\documentclass{article}
\usepackage[UTF8]{ctex}
\usepackage{amsmath,amssymb,amsfonts}
\usepackage{enumitem}
\usepackage{geometry}
\geometry{a4paper, margin=1in}
% === 错题本排版优化（不参与题目标识） ===
\setlength{\parindent}{0pt}
\setlength{\parskip}{0.35em plus 0.10em minus 0.08em}
\setlength{\emergencystretch}{3em}
\setlist{itemsep=0.25em, topsep=0.40em, parsep=0pt}
\allowdisplaybreaks[2]
% === 错题本排版优化结束 ===
\begin{document}
\section*{暨南大学2024年数学分析真题}
\section*{158．暨南大学2024年数学分析试卷}
合计算题．

\subsection{题目1}
1．求极限 \(\lim _{n \rightarrow \infty} \cos \frac{x}{2} \cos \frac{x}{4} \cdots \cos \frac{x}{2^{n}}(x \neq 0)\) ．

\subsection{题目2}
2．求极限 \(\lim _{x \rightarrow 0} \frac{e^{\sin x}-\ln (1+x)-1}{[\ln (1+2 x)]^{2}}\) ．

\subsection{题目3}
3 ．求积分 \(\int_{0}^{2024 \pi} \sqrt{1-\cos 2 x} \mathrm{~d} x\) ．

\subsection{题目4}
4．设 \(y=x^{2} e^{2 x}\) ，求 \(y^{(20)}\) ．

\subsection{题目5}
5．设 \(f(x)=(\arctan x)^{\ln x}\) ，求 \(f^{\prime}(x)\) ．

\subsection{题目6}
6．求不定积分 \(\int \frac{5 x^{2}+1}{x\left(x^{2}+1\right)} \mathrm{d} x\) ．

\subsection{题目7}
7．求定积分 \(\int_{0}^{1} x^{m}(\ln x)^{n} \mathrm{~d} x, m, n \in \mathbb{N}^{+}\)．

\subsection{题目8}
8．求极限 \(\lim _{n \rightarrow \infty}\left(\frac{n}{n^{2}+1^{2}}+\frac{n}{n^{2}+2^{2}}+\cdots+\frac{n}{n^{2}+n^{2}}\right)\) 。

\subsection{题目9}
9．求 \(2 y=x\) 与 \(y=x^{2}\) 围成区域绕 \(y\) 轴旋转所得旋转体的体积。

\subsection{题目10}
1．已知 \(a_{n}, b_{n}\) 是正数列，\(k_{n}=\frac{b_{n} a_{n}}{a_{n+1}}-b_{n+1}\) ，存在 \(c>0\) ，使得对任意的正整数 \(n\) ，均有 \(k_{n}>c\) ，证明：级数 \(\sum_{n=1}^{\infty} a_{n}\) 收敛。

\subsection{题目11}
2．设 \(f_{0}, f_{1}\) 可微，且 \(f_{0} \neq 0,\left|\begin{array}{ll}f_{0} & f_{1} \\ f_{0}^{\prime} & f_{1}^{\prime}\end{array}\right| \neq 0\) ，证明：对任意的 \(\lambda_{0}, \lambda \in \mathbb{R} \backslash\{0\}\) ，函数 \(F(x)=\lambda_{0} f_{0}(x)+ \lambda_{1} f_{1}(x)\) 至多存在一个零点，且若存在零点，则该点 \(F(x)\) 的一阶导数不为零．

\subsection{题目12}
3．设 \(P(x, y), Q(x, y)\) 可微，区域 \(D\) 是由曲线 \(L:\left\{\begin{array}{l}x=\varphi(t) \\ y=\psi(t)\end{array}\right.\) 所围成的区域，且 \(P_{x}+Q_{y}=0\) 无根．证明：

\[
\left\{\begin{array}{l}
\varphi^{\prime}(t)=P(\varphi(t), \psi(t)) \\
\psi^{\prime}(t)=Q(\varphi(t), \psi(t))
\end{array}\right.
\]

一定不成立。

\subsection{题目13}
4．设 \(f(x)\) 在 \([a, b]\) 上连续，且 \(\int_{a}^{b} f^{2}(x) \mathrm{d} x=0\) ，证明：\(f(x)\) 在 \([a, b]\) 上恒为零．

\subsection{题目14}
1．分析 \(f(x)=x^{2} \cos \left(\frac{\pi}{2}-x\right)-x^{3}+x^{4}\) 在 0 处是否取得极值．

\subsection{题目15}
2．设 \(f_{n}(x)=e^{n(x-1)}, x \in(0,1)\) ，分析 \(\left\{f_{n}(x)\right\}\) 在 \((0,1)\) 上的收敛性和一致收敛性．
\end{document}
